\section{Problem Description}\label{sec:problem}

\subsection{Problem statement}

The brief asks for a Java Swing or JavaFX task management (TODO) system in which each task has a title, a description
and a status. We chose JavaFX; the reasons are given in Section~\ref{sec:design}.

\subsection{Functional requirements}

The brief sets six functional requirements, numbered here so that the test cases can refer to them:
\begin{enumerate}
  \renewcommand{\labelenumi}{FR\arabic{enumi}}
  \item add a task, with a title and a description;
  \item view all tasks;
  \item update a task;
  \item delete a task;
  \item set a task's status to Pending, In Progress or Completed;
  \item validate the input, so that no task has an empty title, and handle errors.
\end{enumerate}

\subsection{Requirements we added}

A task list that is lost when the program closes is of little use, and some of the brief's test cases cannot be
shown at all without a little more than the six requirements. We therefore added the following, and nothing beyond
it:
\begin{itemize}
  \item tasks are kept in a database file and survive a restart;
  \item the list can be searched by title or description, and filtered by status;
  \item the number of tasks in each status is always in view;
  \item a task's status can be changed straight from the list, as well as in the edit form;
  \item each task shows when it was created and when it was last updated;
  \item the title and the description have length limits, and both are trimmed of surrounding spaces, so that a
        title of spaces only counts as empty;
  \item deleting a task asks for confirmation first;
  \item adding, editing, deleting, searching and refreshing have keyboard shortcuts;
  \item a task that no longer exists is handled gracefully, as explained below.
\end{itemize}

\subsection{A task that no longer exists}

One test case in the brief needs a word of explanation: updating a non-existing task must be handled gracefully. In a
graphical program the list only offers tasks that exist, so a task can only be missing if something else removed it
after the list was shown. With the tasks in a database file, that something can be a second copy of the program open
on the same file. If one copy deletes a task while the other still shows it, and the second copy then saves a change
to it, deletes it or changes its status, the program must say plainly that the task no longer exists and refresh its
list, without a stack trace. This makes the test case reproducible by hand with two copies of the program open, and
Section~\ref{sec:testing} does exactly that.

\subsection{Constraints}

The program is a single-user desktop application for Windows. It needs no server and no installation beyond the JDK:
its data is a single SQLite file in the signed-in user's local application data folder,
\path{%LOCALAPPDATA%\TaskManager\taskmanager.db}. Every error the user can cause is answered with a message in plain
words; none reaches the user as a stack trace.
